% --- Άσκηση 34317 (34317.tex) ---
\Askhsh{34317}{4}
\begin{ylh}{2}
    \item 4.1. Ανισώσεις 1ου Βαθμού
    \item 6.1. Η Έννοια της Συνάρτησης
\end{ylh}
Το ποσό που θα πληρώσει (σε ευρώ) ένας κάτοικος μιας πόλης $A$ ο οποίος καταναλώνει $x$ κυβικά μέτρα νερού σε ένα χρόνο, δίνεται από τη συνάρτηση:
\[ f(x) = \begin{cases} 0,5x + 12, & \text{αν } 0 \le x \le 30 \\ 0,7x + 6, & \text{αν } x > 30 \end{cases} \]
\begin{alist}
\item Να βρείτε πόσα ευρώ θα πληρώσει κάποιος αν:
\begin{rlist}
\item
  έλλειπε από το σπίτι του και δεν έχει καταναλώσει καθόλου νερό,
\monades{2}
\item
  έχει καταναλώσει $10$ κυβικά μέτρα νερού,
\monades{3}
\item
  έχει καταναλώσει $50$ κυβικά μέτρα νερού.
\monades{5}
\end{rlist}
\item Σε μια άλλη πόλη $B$, το ποσό (σε ευρώ) που αντιστοιχεί σε κατανάλωση $x$ κυβικών μέτρων δίνεται από τον τύπο: $g(x) = 12 + 0,6x,$ για $x \geq 0$.
Ένας κάτοικος της πόλης $A$ και ένας κάτοικος της πόλης $B$ κατανάλωσαν τα ίδια κυβικά μέτρα νερού. Αν ο κάτοικος της πόλης $A$ πλήρωσε μεγαλύτερο ποσό στο λογαριασμό του από τον κάτοικο της πόλης $B$, να αποδείξετε ότι ο κάθε ένας από τους δύο κατανάλωσε περισσότερα από $60$ κυβικά μέτρα νερού.
\monades{15}
\end{alist}
